\section{Executable predictive quotients and the retained flow mechanism}
\label{sec:quotient}
Fix a declared class $\mathcal T$ of executable future tests, including the legal command domains, resource records and scored readouts. Suppose a countable family $(T_j)$ determines its conditional laws. For finite visible histories $h$ set
\[
 q(h)=(\law(T_j\mid h))_{j\ge1}.
\]
Take Borel conditional versions on the histories under consideration. Assume the actual image $S=q(\mathsf H)$ is Borel and admits a Borel section, legal command domains are constant on fibers, and raw report/update laws factor jointly measurably through $q$. These are realization hypotheses; arbitrary measurable equivalence relations need not satisfy them.

\begin{proposition}[Minimality for declared executable tests]\label{prop:quotient}
Under these assumptions, $q$ realizes exactly the predictive equivalence relation for $\mathcal T$ and carries a jointly Borel report/update kernel. Every other Borel statistic that determines these tests and admits the corresponding Borel factorization carries $q$ measurably. For a stopped test, equality of its law follows only with the declared almost-sure termination or approximation conditions.
\end{proposition}
\begin{proof}
Coordinate conditional laws are Borel. Their equality is precisely predictive equivalence by determination. Compose the raw kernel with a Borel section of $q$ and push the continued history through $q$. Fiber constancy makes the result independent of that section and jointly Borel in state and command. Induction proves equality for finite executable prefixes. For bounded stopped readouts under almost-sure termination, truncate at time $n$ with a distinguished surviving symbol and let $n\to\infty$; the surviving probability vanishes. For another sufficient statistic the test coordinates are constant on its fibers and the assumed factorization supplies the Borel map to $q$. This is minimality for the declared tests, not an unrestricted optimal-control or infinite-horizon equivalence.
\end{proof}

The first and second task moments in the main theorem are functions of this predictive information when its realization hypotheses hold, but they need not separate all future tests. A known finite hidden-state or quantum model has its usual conditional-state carrier; if the tests do not determine that carrier, its predictive quotient is smaller. For an unknown fixed model, the family of modelwise posteriors is a comparison object, not a single executable observer. The upper in Theorem~\ref{thm:main} stores neither the exact posterior nor that family.

For comparison with classical occupation methods, we state the inherited submeasure-preserving realization separately. It is not a hidden premise of the new balanced theorem.

\begin{proposition}[Finite occupation on a Borel predictive carrier]\label{prop:flow}
Let $S,A,Y$ be standard Borel with legal Borel domains, a Borel legal default policy on all states, and transition/report kernel $K(ds',dy\mid s,a)$. Let $\alpha$ be a prepared probability measure. Suppose a finite nonnegative legal occupation measure $\eta(ds,da)$ and a legally supported terminal measure $\zeta$ satisfy
\[
 0\le\zeta\le\eta K,\qquad
 \eta_S=\alpha+(\eta K-\zeta)_{s'}.
\]
There is a stationary randomized policy with a Borel post-report stopping rule whose occupation is at most $\eta$ and whose full terminal measure is exactly $\zeta$. It terminates almost surely with finite expected count and no greater nonnegative additive costs.
\end{proposition}
\begin{proof}
Disintegrate $\eta=\xi\pi$ and choose a clipped Borel density $h=d\zeta/d(\eta K)$, zero outside the stopping domain. The continued-state kernel is
$Q(s,D)=\int\pi(da\mid s)\int\one_D(s')(1-h)K(ds',dy\mid s,a)$.
Balance is $\xi=\alpha+\xi Q$; its total mass also implies $\zeta(1)=1$. Monotone iteration gives the actual minimal occupation $\bar\xi=\sum_{j\ge0}\alpha Q^j\le\xi$. The residual $d=\xi-\bar\xi$ is finite and satisfies $d=dQ$, hence $\int(1-Q\one)\,dd=0$. Only rows of zero acceptance carry residual mass. With $r=d\bar\xi/d\xi$, the realized terminal measure is $r(s)\zeta$, and its missing mass is exactly $\int(1-Q\one)\,dd=0$. Thus the terminal measure is unchanged. The actual occupation $\bar\eta=\bar\xi\pi$ is dominated by $\eta$, proving the count and cost bounds. Finite expected count implies almost-sure commitment. Null-set choices can be completed by any measurable legal policy on the unreachable part; existence of such a completion is part of the declared Borel legal interface.
\end{proof}

This is a fixed-kernel Markovization statement adjacent to \cite{PZExtreme,PZContinuity,DPAbsorb}. It neither gives a finite-state implementation of an arbitrary Borel policy nor merges separately chosen unknown-model policies into one program. The common finite-state implementation here comes instead from the raw moment balance and the counted rejection rule.
